\documentclass[10pt,a4paper]{amsart}
\usepackage[margin=19mm]{geometry}
\usepackage{amsmath,amssymb,mathtools}
\usepackage[scr=rsfs,cal=euler]{mathalfa}
\usepackage{xfrac}
\usepackage[unicode,hidelinks,bookmarks=false]{hyperref}
\newcommand{\ie}{i.e.\ }
\newcommand{\cf}{cf.\ }
\newcommand{\resp}{respectively\ }
\newcommand{\Subsection}{\subsection}
\providecommand{\nobreakdash}{\nobreak\hbox{-}}
\setlength{\emergencystretch}{2em}
\multlinegap0pt
\usepackage{enumerate}
\usepackage{amssymb}
\usepackage{tikz}
\usepackage{xcolor}
\newtheorem{proposition}{Proposition}
\newtheorem{lemma}{Lemma}
\newtheorem{corollary}{Corollary}
\def\C{\mathbb{C}}
\def\R{\mathbb{R}}
\newcommand{\be}{\mathbf e}
\newcommand{\bmu}{{\boldsymbol{\mu}}}
\newcommand{\bv}{\mathbf v}
\def\ii{\mathrm{i}}
\newcommand\spec{{\mathrm{spec}}}
\newcommand{\supp}{{\Phi}}
\title{Local $\epsilon$-uniform mixing in continuous quantum walks}
\author{Hermie Monterde}
\date{Source: arXiv:2603.20977v1 (CC BY 4.0)}
\begin{document}
\maketitle

\noindent\emph{Source and licence.} Local $\epsilon$-uniform mixing in continuous quantum walks. Version arXiv:2603.20977v1 (CC BY 4.0), distributed by arXiv under the Creative Commons Attribution 4.0 International licence, http://creativecommons.org/licenses/by/4.0/, as recorded in the arXiv licence metadata for this exact version and shown on the abstract page https://arxiv.org/abs/2603.20977v1. Full licence notice: \texttt{licenses/arxiv-2603.20977-CC-BY-4.0.txt}.\par\smallskip

\noindent\textit{Editorial scope.} Counted unit: the article's corollary nodd (source0.tex lines 882--885). The article's complete original proof of it is delivered with it (source0.tex lines 887--900, one \texttt{proof} environment of 14 lines). No mathematics is written, repaired or re-derived: the statement and the proof are the article's own lines, in the article's own notation, and no retained line is substituted or re-flowed.  Retained uncounted as the source-local context the proof needs: lines 178--194; lines 281--298; lines 317--342; lines 783--790.  Nothing was omitted from any retained span, so the package carries no omission record.  No retained line contains a citation, so the wrapper carries no bibliography and no bibliography entry.  No retained line contains a figure environment, an image-inclusion command or drawing code, so no image is delivered and the package declares no asset dependency.  Editorial formatting only, in addition: an AMS \texttt{amsart} preamble that restates the article's own declarations and macros used by the retained lines, namely the environments \texttt{proposition}, \texttt{lemma}, \texttt{corollary} on separate counters, together with the standard packages those lines need.  The retained environments print in this document as Proposition 1, Lemma 1, Corollary 1 in order of appearance, whereas the article prints the same items with the numbering of its own full document.  The complete native source of the article as distributed in the arXiv e-print for this exact version is bundled unchanged as \texttt{sources/196903/source0.tex}.
\begin{proposition}
\label{proplum}
Let $X$ be a weighted graph on $n$ vertices. The following are equivalent.
\begin{enumerate}
\item Local $\epsilon$-uniform mixing occurs at $u$.
\item There is a sequence $\{\tau_m\}\subseteq\R$ and a 1-uniform vector $\bmu\in\C^n$ such that
\begin{equation}
\label{locum}
\lim_{m\rightarrow\infty}\sqrt{n}U(\tau_m)\be_u=\bmu.
\end{equation}
\item There is a sequence $\{\tau_m\}\subseteq\R$ and a 1-uniform vector $\bmu\in\C^n$ such that 
\begin{center}
$\big(\displaystyle\lim_{m\rightarrow\infty}e^{\ii \tau_m\lambda}\big)\sqrt{n}E_\lambda\be_u=E_\lambda\bmu\quad $ for all $\lambda\in\supp_{\be_u}$.
\end{center}
Equivalently, for each eigenpair $(\bv,\lambda)$ with $\lambda\in\supp_{\be_u}$, we have $\big(\displaystyle\lim_{m\rightarrow\infty}e^{\ii\tau_m\lambda}\big)\sqrt{n}(\bv^T\be_u)=(\bv^T\bmu)$. Thus, $\displaystyle\lim_{m\rightarrow\infty}e^{\ii\tau_m\lambda}$ exists for all  $\lambda\in\supp_{\be_u}$.
\end{enumerate}
\end{proposition}

\begin{lemma}
\label{Lem:no1col}
With the assumption in Proposition~\ref{proplum}, the following hold. 
\begin{enumerate}
\item For each $\lambda\in\supp_{\be_u}$, $\sqrt{n}\left\|E_\lambda\be_u\right\|=\left\|E_\lambda\bmu\right\|$. Consequently, $\supp_{\be_u}=\supp_{\bmu}$.
\item 
For each $\lambda\in\spec(M)$, %$\displaystyle n(E_\lambda)_{u,u}=m_\lambda+2\sum_{j>\ell}\cos(\theta_j-\theta_{\ell})(E_\lambda)_{j,{\ell}},$
\begin{equation}
\label{111}
n(E_\lambda)_{u,u}=m_\lambda+2\sum_{j>\ell}\cos(\theta_j-\theta_{\ell})(E_\lambda)_{j,{\ell}},
\end{equation}
\item For all integers $r\geq 0$,
\begin{equation}
\label{cosp}
n(M^r)_{u,u}=\sum_{j\in V(X)}(M^r)_{j,j}+2\sum_{j>\ell}\cos(\theta_j-\theta_{\ell})(M^r)_{j,{\ell}}.
\end{equation}
\end{enumerate}
\end{lemma}

\begin{corollary}
\label{Cor:eqs}
Let $X$ be an unweighted graph on $n$ vertices. If local $\epsilon$-uniform mixing occurs at $u$ with target state $\bmu=[e^{\ii \theta_1},e^{\ii \theta_2},\ldots,e^{\ii \theta_{n}}]^T$, then the following hold.
\begin{enumerate}
\item If $M=A$, then
\begin{equation}
\label{Eq:cos}
\sum_{\{j,\ell\}\in E(X)}\cos(\theta_j-\theta_{\ell})=0,
\end{equation}
and
\begin{equation}
\label{Eq:cos1}
|E(X)|+\sum_{j>\ell}\cos(\theta_j-\theta_{\ell})c_{j,\ell}=\frac{1}{2}n\operatorname{deg}u.
\end{equation}
\item If $M=Q$, then
\begin{equation}
\label{Eq:cosQ}
|E(X)|+\sum_{\{j,\ell\}\in E(X)}\cos(\theta_j-\theta_{\ell})=\frac{1}{2}n\operatorname{deg}u.
\end{equation}
\item If $M=L$, then
\begin{equation}
\label{Eq:cosL}
|E(X)|-\sum_{\{j,\ell\}\in E(X)}\cos(\theta_j-\theta_{\ell})=\frac{1}{2}n\operatorname{deg}u.
\end{equation}
\end{enumerate}
\end{corollary}

\begin{proposition}
\label{1}
Let $X$ be a weighted bipartite graph with parts $B_1$ and $B_2$.
\begin{enumerate}
\item If local $\epsilon$-uniform mixing occurs at $u\in B_1$ (resp., $u\in B_2$), then the target state is $[\bmu_1,\ii\bmu_2]^T$ (resp., $[\ii\bmu_1,\bmu_2]^T$), where $\bmu_1$ and $\bmu_2$ are $\pm 1$-vectors indexed by $B_1$ and $B_2$, respectively.
\item If $X$ has uniform mixing in $X$ at $\tau$, then $X$ is periodic and $\sqrt{n}U(\tau)$ is a Turyn Hadamard matrix.
\end{enumerate}
\end{proposition}

\begin{corollary}
\label{nodd}
Let $X$ be an unweighted bipartite graph on $n$ vertices admitting local $\epsilon$-uniform mixing at vertex $u$. Then $n$ or $\operatorname{deg}u$ is even. Moreover, there is an even number of vertices in $X$ with degree $d\equiv 2,3$ (mod 4) if and only if $|E(X)|$ and $\frac{1}{2}n\operatorname{deg}u$ have the same parities.
\end{corollary}

\begin{proof}
By Proposition~\ref{1}, the vector $\bmu$ in Lemma~\ref{Lem:no1col} satisfies $\theta_j=\frac{m_j\pi}{2}$ for each $j$, where $m_j$ is some integer. Thus, $\cos(\theta_j-\theta_{\ell})$ is an integer for all $j,\ell\in V(X)$ with $j>\ell$. Invoking (\ref{Eq:cos1}) in Corollary~\ref{Cor:eqs}, we find that $\frac{1}{2}n\operatorname{deg}u$ is an integer. Hence, $n$ or $\operatorname{deg}u$ is even. To complete the proof, we utilize (\ref{Eq:cos1}), which states that
\begin{equation}
\label{eq0}
|E(X)|+\sum_{j>\ell}\cos(\theta_j-\theta_{\ell})c_{j,\ell}=\frac{1}{2}n\operatorname{deg}u.
\end{equation}
If $j$ and $\ell$ belong to different parts, then $c_{j,\ell}=0$ because $X$ is bipartite. On the other hand, if $j$ and $\ell$ belong to the same parts, then $c_{j,\ell}$ may or may not be equal to 0. For the case when $c_{j,\ell}\neq 0$, Proposition~\ref{1}(1) implies that $\cos(\theta_j-\theta_{\ell})=\pm 1$. Thus, $\sum_{j>\ell}\cos(\theta_j-\theta_{\ell})c_{j,\ell}$ has the same parity as $\sum_{j>\ell}c_{j,\ell}$. 
If $v$ is a pendent vertex in $X$, then $v$ is not a common neighbor of any two vertices. Otherwise, $\operatorname{deg}v\geq 2$ and there are ${\operatorname{deg}v\choose 2}$ distinct pairs of vertices having $v$ as a common neighbor. Thus, 
\begin{equation}
\label{q}
\sum_{j>\ell}c_{j,\ell}=\sum_{{v\in V(X)}\atop{\operatorname{deg}v\geq 2}}{\operatorname{deg}v\choose 2},
\end{equation}
where the expression on the right counts the total number of common neighbors of every pair of vertices in $T$. Finally, since $\sum_{j>\ell}\cos(\theta_j-\theta_{\ell})c_{j,\ell}$ and $\sum_{j>\ell}c_{j,\ell}$ have the same parities, and ${\operatorname{deg}v\choose 2}$ is even if and only if $\operatorname{deg}v\equiv 0,1$ (mod 4). Thus, $\sum_{j>\ell}\cos(\theta_j-\theta_{\ell})c_{j,\ell}$ is even if and only if there is an even number of vertices in $X$ with degree $d\equiv 2,3$ (mod 4). Combining this with (\ref{eq0}) yields the desired conclusion.
\end{proof}

\end{document}
